
% Tome v2 – Block B
% Core Genesis Algebra and White State Core

\chapter{Genesis Algebra and the Mote}

In this chapter we introduce the \emph{mote} as the first stable, discrete
geometric excitation above the zero state.  We treat the mote as a
two-layer $2\times 2$ lattice of Planck-scale beads with exactly one
diagonal pair occupied on each layer.  The bottom layer encodes an
``ignition'' choice (left- or right-handed), while the top layer is fixed
by the diagonal rule.

\section{From $0+0$ to the first stable object}

We begin from a purely algebraic statement of nothingness: two
independent null slots,
\[
  (0,0),
\]
with no distinguished ordering and no background manifold.  A \emph{mote
creation} event is defined as the minimal operation that:

\begin{enumerate}
  \item preserves overall null content (no net classical quantity is
  created),
  \item but breaks the degeneracy between the two slots by assigning
  complementary geometric roles.
\end{enumerate}

We represent this by the schema
\begin{equation}
  0 \oplus 0 \;\longrightarrow\; 1,
  \label{eq:zero-plus-zero-equals-one}
\end{equation}
where $\oplus$ is \emph{not} ordinary addition, but a constrained
Genesis operation that produces a single stable mote state from two
indistinguishable null slots.  Equation~\eqref{eq:zero-plus-zero-equals-one}
is the defining axiom of \emph{Genesis algebra}.

\section{Lattice representation}

Let the four sites of a $2\times 2$ layer be labelled
\[
  (x,y) \in \{0,1\}\times\{0,1\},
\]
with $x$ horizontal and $y$ vertical.  A single layer configuration is
then a binary function
\[
  n(x,y) \in \{0,1\},
\]
with the \emph{diagonal rule}
\begin{equation}
  \sum_{x,y} n(x,y) = 2,
  \qquad
  n(0,0) + n(1,1) = 1,
  \qquad
  n(0,1) + n(1,0) = 1.
  \label{eq:diagonal-rule}
\end{equation}
Thus each layer contains exactly one \emph{occupied diagonal}.

The full mote consists of a pair of such layers,
\[
  n_{\mathrm{bot}}(x,y),\quad
  n_{\mathrm{top}}(x,y),
\]
with the constraint that the top diagonal is the complement of the
bottom diagonal.  If the bottom layer is
\[
  n_{\mathrm{bot}}(0,0) = n_{\mathrm{bot}}(1,1) = 1,\quad
  n_{\mathrm{bot}}(0,1) = n_{\mathrm{bot}}(1,0) = 0,
\]
then the top layer is
\[
  n_{\mathrm{top}}(0,1) = n_{\mathrm{top}}(1,0) = 1,\quad
  n_{\mathrm{top}}(0,0) = n_{\mathrm{top}}(1,1) = 0.
\]

\section{Chirality and handedness}

We define a chirality label
\begin{equation}
  \chi = 
  \begin{cases}
    +1, & \text{if the bottom layer occupies the $(0,0)\!\leftrightarrow\!(1,1)$ diagonal},\\[4pt]
    -1, & \text{if the bottom layer occupies the $(0,1)\!\leftrightarrow\!(1,0)$ diagonal}.
  \end{cases}
  \label{eq:mote-chirality}
\end{equation}
The top layer then necessarily has chirality $-\chi$.

At the purely algebraic level, the two chiralities are equivalent; the
Genesis algebra itself does not prefer left or right.  Physical
symmetry-breaking (for example, in the emergence of right-handed
biomolecular structures) can then be modelled as a bias in the
statistical distribution of $\chi$ over large ensembles of motes.

\section{Hilbert space and creation operator}

We now embed the discrete mote configurations in a minimal Hilbert
space.  Let $\mathcal{H}_{\mathrm{mote}}$ be spanned by the two
orthonormal chirality eigenstates
\[
  \ket{+} \equiv \ket{\chi=+1},\qquad
  \ket{-} \equiv \ket{\chi=-1}.
\]
The vacuum state $\ket{\Omega}$ carries no motes:
\[
  \hat{N} \ket{\Omega} = 0,
\]
where $\hat{N}$ is a number operator.

A single \emph{Genesis creation operator} $\hat{G}^\dagger$ is defined by
\begin{equation}
  \hat{G}^\dagger \ket{\Omega}
   \;=\;
   \frac{1}{\sqrt{2}}\bigl(\ket{+} + \ket{-}\bigr),
  \label{eq:genesis-creation}
\end{equation}
creating an equal superposition of left- and right-handed motes from
the null state.  Its adjoint acts as an annihilation operator,
\[
  \hat{G}\ket{+} = \hat{G}\ket{-} = \frac{1}{\sqrt{2}} \ket{\Omega},
\]
and we impose the algebra
\begin{equation}
  \hat{G}\hat{G}^\dagger \ket{\Omega} = \ket{\Omega},\qquad
  \hat{G}^\dagger \hat{G} \ket{\pm} = \ket{\pm},
  \label{eq:genesis-algebra}
\end{equation}
so that $\hat{G}$ and $\hat{G}^\dagger$ act as partial isometries on
the single-mote sector.

\chapter{Infinity Spectrum and Collapse Algebra}

In this chapter we extend Genesis algebra to an \emph{infinity spectrum}
of possible mote excitations, and introduce a simple \emph{collapse
algebra} that models how large families of motes reduce to stable,
finite configurations suitable for building time-shell lattices.

\section{Definition of the infinity spectrum}

We label distinct ``sizes'' or complexity classes of infinity by an
index $k \in \mathbb{N}$ and write the corresponding abstract infinity
as $\mathfrak{I}_k$.  We demand only that
\begin{equation}
  \mathfrak{I}_1 \prec \mathfrak{I}_2 \prec \cdots,
  \label{eq:infinity-spectrum}
\end{equation}
where $\prec$ denotes ``strictly smaller than'' in the sense of
cardinality.  No explicit set-theoretic construction is required at this
stage; we treat $\{\mathfrak{I}_k\}$ as a formal spectrum.

\section{Collapse maps}

A \emph{collapse map} is a homomorphism
\[
  C_k : \mathfrak{I}_k \longrightarrow \mathcal{H}_{\mathrm{mote}},
\]
sending an abstract infinity to a normalized mote superposition.  We
represent this schematically by
\begin{equation}
  C_k(\mathfrak{I}_k)
    \;=\;
    \alpha_k \ket{+} + \beta_k \ket{-},
  \qquad
  |\alpha_k|^2 + |\beta_k|^2 = 1.
  \label{eq:collapse-map}
\end{equation}
Different choices of $k$ correspond to different \emph{collapse
profiles}, e.g. biased vs.\ unbiased chirality outcomes.

\section{White state core as infinity engine}

Finally we combine these ingredients into a minimal phenomenological
Hamiltonian for the white state core,
\begin{equation}
  \hat{H}_{\mathrm{core}}
   \;=\;
   \sum_{k} \lambda_k
      \bigl(
        \hat{G}^\dagger C_k(\mathfrak{I}_k)
        + C_k(\mathfrak{I}_k)^\dagger \hat{G}
      \bigr)
   \;+\; \hat{H}_{\mathrm{self}},
  \label{eq:white-core-hamiltonian}
\end{equation}
where $\lambda_k$ are coupling constants and $\hat{H}_{\mathrm{self}}$
collects self-interaction terms in the mote sector.  The core thus acts
as an \emph{infinity engine}: it continuously converts abstract
infinite degrees of freedom into concrete, chirality-labelled motes
that can be assembled into higher structures such as ER-foam lattices
and time shells.

This block provides the detailed algebraic backbone for the conceptual
discussion in Chapter~0 and prepares the ground for the later,
fully-fledged field-theoretic treatment of helicons, dark bridges, and
cosmological phase structure.
